\setlength{\baselineskip}{20pt}
\begin{abstract}
去噪模型作为扩散框架的核心机制，已在生成建模中展现出强大的分布刻画与样本合成能力，但其在判别任务中的潜力仍有待系统挖掘。判别学习本质上依赖表征学习，即学习能够支撑分类与预测的判别性特征；然而真实视觉数据普遍伴随噪声污染、域偏移与背景干扰，容易造成类内离散增大和类间边界模糊。已有扩散思想在判别场景中的应用多针对单一任务，通用性不足，且常需显式递归去噪步骤，导致推理延迟与部署成本上升。围绕“如何在保持计算效率的同时提升特征判别性”这一关键问题，本文以“特征提取—特征去噪”统一建模为主线展开研究。本文的研究内容与主要贡献可归纳为以下两个方面：
\par
（1）面向无额外推理成本的特征增强算法研究：将判别骨干网络中的级联嵌入层解释为逐层去噪过程，实现特征提取与特征去噪的统一建模；进一步构建去噪层与特征层的参数融合机制，给出融合前后严格等价的变换关系，将训练阶段的双分支增强折叠为推理阶段的单分支计算，从而在不改变任务头和参数规模的前提下实现无额外推理成本的特征增强；同时兼容无标签训练与有标签监督训练，使去噪信号与判别信号能够协同优化，并可直接嵌入多类判别框架，迁移成本较低且部署流程简洁。
\par
（2）基于模型蒸馏的特征增强算法研究：针对轻量去噪模块表达能力受限、噪声估计偏差较大的问题，引入教师--学生蒸馏机制，以高容量教师网络提供更稳定的噪声方向监督；提出多轮单层去噪与多步去噪融合策略，在不增加部署复杂度的条件下压缩多步轨迹并增强有限去噪深度下的表征能力；通过消融与对比评估验证蒸馏与融合模块的互补性与稳定收益，使方法在同等预算下获得更强的跨任务泛化性能。进一步地，蒸馏目标与多步融合可统一训练并稳定收敛。
\par
实验覆盖行人重识别、车辆重识别、图像分类、目标检测与语义分割等任务，包含 CNN 与 Transformer 等多类骨干网络以及多个公开数据集，评价指标包括 mAP/Rank-1、Top-1/Top-5、AP 系列指标与 mIoU。结果表明，所提方法在不同骨干与不同任务上均可获得稳定性能增益，并在计算预算受限条件下保持较高效率，验证了方法的工程可部署价值，为扩散思想服务判别式表征学习提供了统一、有效且可落地的技术路径；同时，两类方法可作为可插拔增强模块适配不同判别骨干，为后续面向更复杂开放场景的鲁棒表征学习提供了可复用的技术基础。相关分析进一步表明，该框架在不同噪声强度、训练规模与数据扰动条件下仍保持稳定收益，体现出良好的跨场景适应能力与工程迁移潜力。
\par\noindent\keywords{扩散模型；判别式表征学习；特征增强；参数融合；模型蒸馏}
\end{abstract}
